\documentclass[10pt,a4paper]{amsart}
\usepackage{amsmath,amssymb,amsthm,graphicx}
% Retained minimal preamble for the Stacks Project chapter flat.tex (More on Flatness).
% Verbatim excerpt of the chapter preamble: Xy-pic setup (native preamble lines 14-18) and the native
% theorem declarations with the mathematical macros (native preamble lines 154-217). Copied mechanically;
% nothing added, removed or redefined. The wrapper class and the amsmath/amssymb/amsthm/graphicx base
% packages are supplied by the assembler, so the class/externaldocument part of the native preamble is omitted.
% For commutative diagrams we use Xy-pic
\usepackage[all]{xy}

% We use 2cell for 2-commutative diagrams.
\xyoption{2cell}
% Theorem environments.
%
\theoremstyle{plain}
\newtheorem{theorem}[subsection]{Theorem}
\newtheorem{proposition}[subsection]{Proposition}
\newtheorem{lemma}[subsection]{Lemma}

\theoremstyle{definition}
\newtheorem{definition}[subsection]{Definition}
\newtheorem{example}[subsection]{Example}
\newtheorem{exercise}[subsection]{Exercise}
\newtheorem{situation}[subsection]{Situation}

\theoremstyle{remark}
\newtheorem{remark}[subsection]{Remark}
\newtheorem{remarks}[subsection]{Remarks}

\numberwithin{equation}{subsection}

% Macros
%
\def\lim{\mathop{\mathrm{lim}}\nolimits}
\def\colim{\mathop{\mathrm{colim}}\nolimits}
\def\Spec{\mathop{\mathrm{Spec}}}
\def\Hom{\mathop{\mathrm{Hom}}\nolimits}
\def\Ext{\mathop{\mathrm{Ext}}\nolimits}
\def\SheafHom{\mathop{\mathcal{H}\!\mathit{om}}\nolimits}
\def\SheafExt{\mathop{\mathcal{E}\!\mathit{xt}}\nolimits}
\def\Sch{\mathit{Sch}}
\def\Mor{\mathop{\mathrm{Mor}}\nolimits}
\def\Ob{\mathop{\mathrm{Ob}}\nolimits}
\def\Sh{\mathop{\mathit{Sh}}\nolimits}
\def\NL{\mathop{N\!L}\nolimits}
\def\CH{\mathop{\mathrm{CH}}\nolimits}
\def\proetale{{pro\text{-}\acute{e}tale}}
\def\etale{{\acute{e}tale}}
\def\QCoh{\mathit{QCoh}}
\def\Ker{\mathop{\mathrm{Ker}}}
\def\Im{\mathop{\mathrm{Im}}}
\def\Coker{\mathop{\mathrm{Coker}}}
\def\Coim{\mathop{\mathrm{Coim}}}

% Boxtimes
%
\DeclareMathSymbol{\boxtimes}{\mathbin}{AMSa}{"02}

%
% Macros for moduli stacks/spaces
%
\def\QCohstack{\mathcal{QC}\!\mathit{oh}}
\def\Cohstack{\mathcal{C}\!\mathit{oh}}
\def\Spacesstack{\mathcal{S}\!\mathit{paces}}
\def\Quotfunctor{\mathrm{Quot}}
\def\Hilbfunctor{\mathrm{Hilb}}
\def\Curvesstack{\mathcal{C}\!\mathit{urves}}
\def\Polarizedstack{\mathcal{P}\!\mathit{olarized}}
\def\Complexesstack{\mathcal{C}\!\mathit{omplexes}}
% \Pic is the operator that assigns to X its picard group, usage \Pic(X)
% \Picardstack_{X/B} denotes the Picard stack of X over B
% \Picardfunctor_{X/B} denotes the Picard functor of X over B
\def\Pic{\mathop{\mathrm{Pic}}\nolimits}
\def\Picardstack{\mathcal{P}\!\mathit{ic}}
\def\Picardfunctor{\mathrm{Pic}}
\def\Deformationcategory{\mathcal{D}\!\mathit{ef}}
\usepackage{cleveref}
\crefname{originaldefinition}{original definition}{original definitions}
\Crefname{originaldefinition}{Original definition}{Original definitions}
\makeatletter
\newlabel{definition-complete-devissage}{{05HH}{source}{}{sourcepointer-1}{}}
\newlabel{definition-complete-devissage@cref}{{[originaldefinition][0][]05HH}{[1][0][]source}{}{}{}}
\makeatother
\crefname{originallemma}{original lemma}{original lemmas}
\Crefname{originallemma}{Original lemma}{Original lemmas}
\makeatletter
\newlabel{lemma-base-change-one-step}{{05H7}{source}{}{sourcepointer-2}{}}
\newlabel{lemma-base-change-one-step@cref}{{[originallemma][0][]05H7}{[1][0][]source}{}{}{}}
\makeatother
\crefname{originallemma}{original lemma}{original lemmas}
\Crefname{originallemma}{Original lemma}{Original lemmas}
\makeatletter
\newlabel{coherent-lemma-affine-base-change}{{02KG}{source}{}{sourcepointer-3}{}}
\newlabel{coherent-lemma-affine-base-change@cref}{{[originallemma][0][]02KG}{[1][0][]source}{}{}{}}
\makeatother
\title{A complete devissage persists under an affine base change}
\author{The Stacks Project Authors}
\date{stacks-project:flat.tex:local-snapshot-20261011; GNU FDL 1.2 or later}
\begin{document}
\maketitle
\noindent\textit{Editorial note.} All delivered mathematical text is the author's own, copied from the pinned Stacks Project chapter file flat.tex; no statement or proof step was rewritten, reordered or removed.\par\medskip
\noindent\textit{Editorial note.} The preamble is the chapter's own retained minimal preamble (sources/stacks-flat/retained-preamble.tex), a verbatim excerpt of the native Xy-pic setup and of the native theorem declarations and mathematical macros; no macro was invented or redefined.\par\medskip
\noindent\textit{Editorial note.} Same-chapter referenced results are located at their \textbackslash{}label line inside the pinned chapter file; the cross-chapter affine base change lemma is located in the project's own tag list staged in the snapshot (sources/stacks-chow/tags.txt), at the verbatim line carrying its canonical tag.\par\medskip
\noindent\textit{Allowed standard prerequisites (editorial).}\begin{itemize}
\item Complete devissage of a finite type module. Same chapter, Definition environment labelled definition-complete-devissage, source line 1195 (Stacks tag 05HH). Supplies the tuple (Z\_k, Y\_k, i\_k, pi\_k, G\_k, alpha\_k) and condition (2) of the definition, which is exactly what the counted proof has to preserve under base change.
\item Base change of a one step devissage. Same chapter, Lemma labelled lemma-base-change-one-step, source line 879 (Stacks tag 05H7). Supplies the one-step half of the conclusion, quoted at proof lines 1288-1290; the remaining cokernel and condition-(2) arguments stay inside the counted span.
\end{itemize}\medskip
\ifdefined\hypertarget\hypertarget{sourcepointer-1}{}\fi
\noindent\textit{Source reference.} Original definition 05HH: sources/stacks-flat/pinned-original-source.tex, source lines 1195--1195. Same-chapter pointer to the definition quoted both in the statement (line 1273) and in the proof (line 1294) for condition (2); the verification that this condition survives base change is the counted obligation and stays in the delivered body.\par\medskip
\ifdefined\hypertarget\hypertarget{sourcepointer-2}{}\fi
\noindent\textit{Source reference.} Original lemma 05H7: sources/stacks-flat/pinned-original-source.tex, source lines 879--879. Same-chapter pointer used at proof line 1289 for the one-step half of the conclusion; the cokernel and condition-(2) arguments that complete the proof are retained verbatim in the counted span.\par\medskip
\ifdefined\hypertarget\hypertarget{sourcepointer-3}{}\fi
\noindent\textit{Source reference.} Original lemma 02KG: sources/stacks-chow/tags.txt, source lines 3371--3371. Cross-chapter standard prerequisite invoked once at proof line 1298 to identify pi'\_\{k,*\}G'\_k with the pullback of pi\_\{k,*\}G\_k; it supplies that identification only, while the cokernel computation and the condition-(2) computation are retained in the counted span.\par\medskip
\begin{lemma}
\label{lemma-base-change-complete}
Let $S$, $X$, $\mathcal{F}$, $s$ be as in
Definition \ref{definition-complete-devissage}.
Let $(S', s') \to (S, s)$ be any morphism of pointed schemes.
Let $(Z_k, Y_k, i_k, \pi_k, \mathcal{G}_k, \alpha_k)_{k = 1, \ldots, n}$
be a complete d\'evissage of $\mathcal{F}/X/S$ over $s$.
Given this data let $X', Z'_k, Y'_k, i'_k, \pi'_k$ be the base
changes of $X, Z_k, Y_k, i_k, \pi_k$ via $S' \to S$.
Let $\mathcal{F}'$ be the pullback of $\mathcal{F}$ to $X'$
and let $\mathcal{G}'_k$ be the pullback of $\mathcal{G}_k$ to $Z'_k$.
Let $\alpha'_k$ be the pullback of $\alpha_k$ to $Y'_k$.
If $S'$ is affine, then
$(Z'_k, Y'_k, i'_k, \pi'_k, \mathcal{G}'_k, \alpha'_k)_{k = 1, \ldots, n}$
is a complete d\'evissage of $\mathcal{F}'/X'/S'$ over $s'$.
\end{lemma}

\begin{proof}
By
Lemma \ref{lemma-base-change-one-step}
we know that the base change of a one step d\'evissage is a one step
d\'evissage. Hence it suffices to prove that formation of
$\Coker(\alpha_k)$ commutes with base change and that
condition (2) of
Definition \ref{definition-complete-devissage}
is preserved by base change. The first is true as
$\pi'_{k, *}\mathcal{G}'_k$ is the pullback of
$\pi_{k, *}\mathcal{G}_k$ (by
Cohomology of Schemes, Lemma \ref{coherent-lemma-affine-base-change})
and because $\otimes$ is right exact. The second because
by the same token we have
$$
(\pi_{k, *}\mathcal{G}_k)_{\xi_k}
\otimes_{\mathcal{O}_{Y_k, \xi_k}} \kappa(\xi_k)
\otimes_{\kappa(\xi_k)} \kappa(\xi'_k)
\cong
(\pi'_{k, *}\mathcal{G}'_k)_{\xi'_k}
\otimes_{\mathcal{O}_{Y'_k, \xi'_k}} \kappa(\xi'_k)
$$
with obvious notation.
\end{proof}

\end{document}
